\subsection{A two-object orbit obstruction in dimension ninety-six}
\begin{proposition}\label{prop:d96-pair}
A fusion category of type $(1^4,2,4,6^2)$ has a nontrivial universal
grading. In particular $R8\_2\_404$, $R8\_2\_418$, $R8\_2\_604$, $R8\_2\_692$
have no complex fusion categorification.
\end{proposition}
\begin{proof}
Suppose such a category has trivial universal grading. Its dimension
is ninety-six, and its invertible group has order four. Solvability
forces a last $C_2$
 equivariantization $\mathcal C=\mathcal D^{C_2}$. Whichever nontrivial
character occurs, it fixes the unique dimension-two and dimension-four
simples. Thus the outer involution exchanges the two dimension-two
simples of $\mathcal D$. Depending on whether it fixes or exchanges
the dimension-six pair upstairs, $\mathcal D$ has type
$(1^4,2^2,3^4)$ or $(1^4,2^2,6)$ and dimension forty-eight.

Its canonical adjoint is preserved by the outer involution.
A proper adjoint's dimension therefore has the form
$a+4b+9c+36e$, where $a\in\{1,2,4\}$, $b\in\{0,2\}$,
$c$ is even, and $e$ is zero or one according to the displayed type.
The proper divisors of forty-eight among these possibilities are only
$1,2,4,12$. A pointed adjoint of dimension at most four cannot contain
the square of a dimension-three or dimension-six simple, since each
invertible occurs there with coefficient at most one. An adjoint of
dimension twelve consists of all four invertibles and both
dimension-two simples. Every other grading component would then have
dimension twelve using only dimension-three or dimension-six simples;
a sum of their squared dimensions is divisible by nine, so this is
impossible. Hence $\mathcal D_{\rm ad}=\mathcal D$.

A second prime descent gives $\mathcal D=\mathcal E^{C_2}$,
$\FPdim(\mathcal E)=24$. The complete cyclic-orbit dimension
partitions are $(1^6,3^2)$ and $(1^2,2,3^2)$.
In the first case the dimension-two pair of $\mathcal D$ comes from
two free orbits in the invertible group $A$ of order six.
The inner involution fixes exactly two elements of $A$.
For $A=C_6$ it is inversion, and for $A=S_3$ it is conjugation by a
transposition. In either group the coefficients of the two induced
objects in their own squares are respectively zero and one, by
ordinary orbit addition. The output stabilizers are trivial, so
associator or tensorator phases cannot change these coefficients.
The outer involution cannot exchange this pair as required.

In the second case the unique dimension-two simple $x$ of
$\mathcal E$ is fixed by the inner involution, with
$x^2=1+g+x$. Its two equivariant lifts $x_+,x_-$ are exchanged by the
central character $\varepsilon$. Therefore their squares are equal,
because $\varepsilon^2=1$ and $\varepsilon$ is central.
Each square has precisely one constituent from the dimension-two pair,
by forgetting to the single $x$ channel. The outer involution exchanges
$x_+,x_-$ but preserves their common square. Its unique dimension-two
constituent would have to be fixed, which is impossible.
Both dimension partitions and every initial character choice have
therefore been excluded. The supplementary check binds the full four
source tensors and exhausts these small dimension and orbit cases.
\end{proof}
